\documentclass[11pt]{article}
\usepackage[T1]{fontenc}
\usepackage{lmodern,microtype}
\usepackage[margin=1in]{geometry}
\usepackage{amsmath,amssymb,amsthm,mathtools,mathrsfs}
\usepackage{booktabs,tabularx,array,longtable}
\usepackage[dvipsnames]{xcolor}
\usepackage[colorlinks=true,linkcolor=MidnightBlue,citecolor=MidnightBlue,urlcolor=MidnightBlue]{hyperref}
\usepackage{bookmark}
\usepackage{enumitem}
\setlist{nosep}
\setcounter{tocdepth}{1}
\renewcommand{\arraystretch}{1.12}
\setlength{\emergencystretch}{2em}
\newtheorem{theorem}{Theorem}[section]
\newtheorem{proposition}[theorem]{Proposition}
\newtheorem{lemma}[theorem]{Lemma}
\newtheorem{corollary}[theorem]{Corollary}
\theoremstyle{definition}
\newtheorem{definition}[theorem]{Definition}
\theoremstyle{remark}
\newtheorem{remark}[theorem]{Remark}
\numberwithin{equation}{section}
\DeclareMathOperator{\Tr}{Tr}
\DeclareMathOperator{\Dom}{Dom}
\DeclareMathOperator{\HS}{HS}
\DeclareMathOperator{\Scal}{Scal}
\newcommand{\R}{\mathbb R}
\newcommand{\C}{\mathbb C}
\newcommand{\K}{\mathcal K}
\newcommand{\HH}{\mathcal H}
\newcommand{\Sch}{\mathcal S}
\newcommand{\dd}{\,\mathrm d}
\newcommand{\ee}{\mathrm e}
\newcommand{\ii}{\mathrm i}
\newcommand{\norm}[1]{\left\lVert#1\right\rVert}
\newcommand{\abs}[1]{\left|#1\right|}
\newcommand{\inner}[2]{\left\langle#1,#2\right\rangle}
\newcommand{\Lsc}{L_{\mathrm{sc}}}
\newcommand{\Dc}{\mathcal D_C}
\newcommand{\Bcal}{\mathcal B}
\newcommand{\Obs}{\mathcal O}
\newcolumntype{Y}{>{\raggedright\arraybackslash}X}
\DeclareMathOperator{\Cov}{Cov}
\DeclareMathOperator{\Var}{Var}
\newcommand{\E}{\mathbb E}
\title{Conditional Newman Dynamics on the Spherical Cylinder\\[4pt]
\large Complete count instruments, calibrated observations,\\
and the anatomy of zero collisions}
\author{Drafted for Edward Baker}
\date{6 November 2026\\\small Working manuscript; analytical results and outstanding tasks}
\hypersetup{pdftitle={Conditional Newman Dynamics on the Spherical Cylinder},pdfauthor={Drafted for Edward Baker},pdfsubject={Direct Newman realization and analytical collision diagnostics}}
\begin{document}
\maketitle
\begin{abstract}
We develop a conditional realization of the canonical de Bruijn--Newman
family on the complete spherical cylinder. Starting from a supplied positive
theta source, we prove regularity of its normalized square root and control
all finite radial and angular energy moments. An explicitly defined count
instrument follows the freely transported position quadrature. Its complete
outcome sum preserves trace, while its no-count branch implements decreasing
Newman time. A bounded cosine observation, together with the success
frequency and one initial calibration, recovers the absolute Newman function.
We then derive the signed moment calculus of every count sector, its generating
identity and transfer laws. These formulas describe a hypothetical zero
collision as a degeneracy of a smooth signed observation: the conditional
state remains regular, while zero tracking has a square-root singularity.
The first count sector carries the curvature at a double event; higher
multiplicities require higher sectors or spatial derivatives. An explicit
smooth reflected control realizes the same generic dynamics and collides,
showing precisely why state positivity and count completeness do not provide
collision control. All proofs are analytical, retain complete spectra and
tails, and distinguish finite-order estimates from uniform claims. The
source loading and apparatus access are supplied assumptions. The final
section identifies the source-specific structural and quantitative tasks
that remain open.
\end{abstract}
\tableofcontents
\clearpage
\section{The direct realization and its analytical purpose}\label{sec:intro}

The direct problem is to realize the Newman deformation through a specified
preparation, evolution and signed observation, and then to understand what
that realization says about zero collisions. The shared manuscript
\cite{TNDBackground} supplies the geometric and spectral background. Here we
retain only the background needed for this branch and derive its preparation,
complete instrument, calibrated observation and collision diagnostics in one
continuous development.

The canonical family is
\begin{align}
 \Phi(v)&=\sum_{n\ge1}(2\pi^2n^4\ee^{9v}-3\pi n^2\ee^{5v})
                       \ee^{-\pi n^2\ee^{4v}},\label{eq:Phi}\\
 H_\lambda(z)&=\int_0^\infty\ee^{\lambda v^2}\Phi(v)\cos(zv)\dd v.
 \label{eq:H}
\end{align}
The complete source is smooth, positive and even; its reflected interpretation
and all limits are justified below. We keep the classical normalization
\begin{equation}\label{eq:xi}
 H_0(z)=\frac18\xi\!\left(\frac{1+\ii z}{2}\right),\qquad
 \xi(s)=\frac12s(s-1)\pi^{-s/2}\Gamma(s/2)\zeta(s),
\end{equation}
as in \cite[Introduction, (1)--(4)]{RodgersTao2020}. No theorem about the
zeros of $H_0$ is used. Direct differentiation gives
$\partial_\lambda H_\lambda=-\partial_z^2H_\lambda$.

The model keeps the free kinetic Hamiltonian on
$\mathcal C=\R\times S^2_{1/2}$. It adds one prescribed count instrument
that monitors the freely transported position quadrature. Its no-count
branch attenuates the positive square-root source by
$\ee^{-\gamma tX^2/2}$ in the interaction frame. The resulting density has
parameter $\lambda(t)=1/2-\gamma t$ on $0\le t\le1/\gamma$.
Physical time therefore runs toward decreasing Newman parameter.

The principal identities organize the paper. If $p_t$ is the no-count
probability, $C(t,k)$ the conditional cosine expectation and $q_\pm$ the joint
success/outcome probabilities, then
\begin{equation}\label{eq:mainintro}
 H_{\lambda(t)}(k)=H_{1/2}(0)\,p_t C(t,k)
                 =H_{1/2}(0)[q_+(t,k)-q_-(t,k)].
\end{equation}
The full instrument has outputs $\sigma_j(t)$, one for each count $j$.
Their signed traces $d_j(t,k)=\Tr[\sigma_j(t)\cos(kQ_t)]$ satisfy
\begin{equation}\label{eq:countintro}
 d_j=\frac{(-\gamma t)^j}{j!}\partial_k^{2j}d_0,\qquad
 \partial_td_0=-d_1/t\quad(t>0),\qquad
 \sum_{j\ge0}d_j=C(0,k).
\end{equation}
Thus the complete count record organizes weighted source moments, while the
Newman evolution appears in its selected no-count sector.

At a hypothetical nondegenerate double event,
$d_0=\partial_kd_0=0$ and $a_*=\partial_k^2d_0\ne0$. The local law is
\begin{equation}\label{eq:foldintro}
 d_0(t_*+h,k_*+x)=\frac{a_*}{2}(x^2+2\gamma h)
                      +O(|x|^3+|hx|+h^2).
\end{equation}
The positive state and success mass remain regular. Equal no-count outcome
probabilities have unequal rates of decrease, encoded by
$d_1(t_*,k_*)=-\gamma t_*a_*$. This provides a dynamical description of a
collision, without asserting that the canonical family has one.

\subsection{Assumptions, evidence and organization}
The source loading, its identification with physical position, the complete
angular preparation, the initial scale $H_{1/2}(0)$ and access to the
specified no-count cosine measurement are supplied model data. The proofs
establish the consistency and exact observation of that model. They do not
derive a microscopic apparatus or an intrinsic theta loader.

The complete realization \cite{TNDDirect} has an internal assessment
accepting its conditional scope \cite{TNDDirectReview}; the shared background
has its own internal assessment \cite{TNDSharedReview}. The later count and
collision analysis \cite{TNDDirectCollision} retains proof-author diagnostic
status. The exposition below makes its derivations explicit but is not an
independent reassessment. No numerical evidence is needed for any proof.

Section~\ref{sec:preparation} establishes the domains and admissible prepared
state. Sections~\ref{sec:instrument}--\ref{sec:observation} construct the
instrument and prove the absolute readout. Section~\ref{sec:count} derives
the complete count calculus. Sections~\ref{sec:collisions}--\ref{sec:control}
analyze local collisions and an exact colliding control. The final section
states the outstanding tasks. The distinction between a representation and
a collision-preventing mechanism remains visible throughout.

\section{Geometric data and admissible preparation}\label{sec:preparation}

\subsection{Complete cylinder and maximal domains}
Let $\K=L^2(S^2,\Sigma;\dd\Omega)$ and
$\HH=L^2(\R,\dd v;\K)$. Let $D$ be the self-adjoint unit-sphere Dirac
operator, $\Gamma$ its chirality, and $A=D^2$. The classical angular spectrum
is $D=\pm n$, $n\ge1$, with multiplicity $2n$ per sign
\cite[(3.40), (3.54)]{CamporesiHiguchi1996}. Write $P_n$ for the rank-$4n$
projection of $A$ at $n^2$. With $P=-\ii\partial_v$ and $X=v$, put
\begin{equation}\label{eq:operators}
 \Dc=\Gamma P+2D,\qquad L_0=\Dc^2=P^2+4A,\qquad
 H_{\rm kin}=\varepsilon L_0,\quad
 \varepsilon=\frac{\hbar^2}{2m_{\rm phys}\ell_0^2},\quad
 \omega=\varepsilon/\hbar>0.
\end{equation}
The product connection and $\Gamma D=-D\Gamma$ remove mixed terms. On
$\dd v^2+g_{S^2}/4$, the scalar curvature is $8$; the Dirac square is
$(\nabla^{\mathcal C})^*\nabla^{\mathcal C}+2$. Thus the fixed curvature
contribution $2\varepsilon$ is already in \eqref{eq:operators}
\cite{TNDBackground,TNDParent}. No potential, wall or radial confinement
is inserted. Metric volume is $\dd v\dd\Omega/4$; multiplication by $2$
identifies product-measure spinors with that convention.

Use $\widehat\psi(k)=(2\pi)^{-1/2}\int\ee^{-\ii kv}\psi(v)\dd v$.
Choosing $\chi^-_{n,a}=\Gamma\chi^+_{n,a}$ makes each Fourier block
\begin{equation}\label{eq:blocks}
 M_n(k)=\begin{pmatrix}2n&k\\k&-2n\end{pmatrix},\quad
 M_n(k)^2=(k^2+4n^2)I,\quad
 \Pi_\pm(n,k)=\frac12\left(I\pm\frac{M_n(k)}{\sqrt{k^2+4n^2}}\right).
\end{equation}
There are $2n$ copies of each sign, all $n\ge1$, and all real $k$.
The square has multiplicity $4n$. The maximal domains are
\begin{align}
 \Dom L_0^p&=\left\{\psi\in\HH:
  \sum_n\int_\R(k^2+4n^2)^{2p}\norm{P_n\widehat\psi(k)}^2\dd k<\infty\right\},
 \label{eq:graphdomain}\\
 \Dom\Dc&=\left\{\psi\in\HH:
  \sum_n\int_\R(k^2+4n^2)\norm{P_n\widehat\psi(k)}^2\dd k<\infty\right\},
 \label{eq:Diracdomain}\\
 \Dom X^q&=\{\psi:|v|^q\psi\in L^2\},\qquad
 \Dom P^q=\{\psi:|k|^q\widehat\psi\in L^2\}.
 \label{eq:XPdomain}
\end{align}
Real Hermitian multiplication on these maximal domains gives the
self-adjoint realizations. Complete-shell cutoffs converge in graph norm;
for finitely many shells the graph norm is a finite sum of Sobolev norms,
so radial cutoff and mollification give a smooth core. In particular
$\Dom L_0=\Dom P^2\cap\Dom A$, because the nonnegative commuting
multipliers dominate their constituents and are bounded by their sum.
The radial spectrum is continuous. Generalized Fourier modes are not
normalizable states, and evenness reduces the square rather than $\Dc$.

\subsection{The source supplied by the background dictionary}
The full-line inverse of $-\partial_v^2+4n^2$ has kernel
$\ee^{-2n|v-w|}/(4n)$: its derivative has jump $-1$.
Its diagonal angular operator is $W=(4\sqrt A)^{-1}$, and
\begin{equation}\label{eq:Q}
 Q(\beta):=\Tr_\K(W\ee^{-\beta A})
 =\sum_{n\ge1}4n\frac{\ee^{-\beta n^2}}{4n}
 =\sum_{n\ge1}\ee^{-\beta n^2},\qquad\beta>0.
\end{equation}
All fixed $\beta$ derivatives converge locally uniformly. The prescribed
scan $\beta(v)=\pi\ee^{4v}$ and $b(v)=\ee^vQ(\beta(v))$ satisfy
\begin{equation}\label{eq:sourceconstruction}
 \Phi(v)=\frac18(b''(v)-b(v)).
\end{equation}
Differentiating the complete series yields \eqref{eq:Phi}. The Green diagonal
is position independent; the scan and imprint are supplied, not consequences
of translating a probe in the free geometry.

Let $\Theta(R)=1+2Q(\pi R^2)$. Jacobi inversion
$\Theta(R)=R^{-1}\Theta(1/R)$ follows from the principal-root specialization
of \cite[20.7.32]{DLMFThetaModular}. With $R=\ee^{2v}$ it gives
$b(v)-b(-v)=-\sinh v$. Applying $(\partial_v^2-1)/8$ proves smooth evenness
of the complete source. For $v\ge0$ every term of \eqref{eq:Phi} is positive,
since $\pi n^2\ee^{4v}>3/2$. Therefore $\Phi>0$ on the full line.
The zero mode in $\Theta$ is essential; finite reflected prefixes cannot
replace this smooth reflection argument.

\begin{lemma}[Complete source estimates]\label{lem:source}
Set $r=|v|$, $x=\pi\ee^{4r}$ and $\kappa=2-3/\pi>0$. For each fixed
$j\ge0$ there is a finite $B_j$ such that
\begin{equation}\label{eq:sourcebounds}
 \kappa\ee^r x^2\ee^{-x}\le\Phi(v),\qquad
 |\Phi^{(j)}(v)|\le B_j\ee^r x^{j+2}\ee^{-x}.
\end{equation}
All polynomial and compact-parameter exponential weights used in this
paper are integrable against these bounds, with complete source tails.
\end{lemma}
\begin{proof}
On $v\ge0$ define
\[
 \mathsf P_0(x)=2x^2-3x,\qquad
 \mathsf P_{j+1}(x)=(1-4x)\mathsf P_j(x)+4x\mathsf P_j'(x).
\]
Differentiation gives
$\Phi^{(j)}(v)=\ee^v\sum_{n\ge1}\ee^{-n^2x}\mathsf P_j(n^2x)$,
with $\deg\mathsf P_j=j+2$. If $a_j$ is its absolute coefficient sum, put
$C_d(N)=\sum_{n>N}n^{2d}\ee^{-\pi(n^2-1)}$. Then
\begin{equation}\label{eq:sourcetail}
 \ee^v\sum_{n>N}\ee^{-n^2x}|\mathsf P_j(n^2x)|
 \le a_jC_{j+2}(N)\ee^v x^{j+2}\ee^{-x},\qquad C_d(N)\longrightarrow0.
\end{equation}
This follows from $x\ge\pi>1$ and Gaussian summability. Taking $N=0$
gives $B_j=a_jC_{j+2}(0)$. The first source term gives the lower bound;
complete reflection transfers both bounds to $v<0$.
For finite $M,B\ge0$ and fixed $a,d\ge0$, substitution
$x=\pi\ee^{4|v|}$ shows that
\begin{equation}\label{eq:J}
 J_{a,d}(M,B)=\int_\R |v|^a\ee^{Mv^2+(B+1)|v|}x^d\ee^{-x}\dd v<\infty.
\end{equation}
The transformed integrand is a fixed power of $x$, powers of $\log(x/\pi)$,
and $\exp[M\log^2(x/\pi)/16-x]$. The term $-x$ dominates all the positive
terms. The same bound with $C_{j+2}(N)$ pays each integrated source tail.
\end{proof}

It follows, by dominated integration on complex parameter compacts, that
$H_\lambda(z)$ is jointly entire and
\begin{equation}\label{eq:Hmixed}
 \partial_\lambda^a\partial_z^bH_\lambda(z)
 =\frac12\int_\R v^{2a}(\ii v)^b\ee^{\lambda v^2}\Phi(v)\ee^{\ii zv}\dd v.
\end{equation}
This proves the Newman equation and all its fixed finite derivatives.
Complex $z$ is an analytic continuation, not a bounded measurement label.

\subsection{The normalized square-root state}
Put $F_\lambda(v)=\ee^{\lambda v^2}\Phi(v)$ and
\begin{equation}\label{eq:f}
 Z_\lambda=\int_\R F_\lambda(v)\dd v=2H_\lambda(0),\qquad
 f_\lambda=Z_\lambda^{-1/2}\ee^{\lambda v^2/2}\sqrt\Phi.
\end{equation}
On $|\lambda|\le M$, $0<Z_{-M}\le Z_\lambda\le Z_M<\infty$.
For instance an explicit lower bound is the integral of the lower estimate
in \eqref{eq:sourcebounds} over $0\le v\le1$, multiplied by
$2\ee^{-Mv^2}$. Define
\begin{equation}\label{eq:R}
 R=\frac{W\ee^{-\pi A}}{Q(\pi)}
   =\sum_{n\ge1}\frac{\ee^{-\pi n^2}}{4nQ(\pi)}P_n,\qquad
 \alpha_b:=\Tr(A^bR)=\frac{\sum_{n\ge1}n^{2b}\ee^{-\pi n^2}}{Q(\pi)}.
\end{equation}
Here $R$ is an angular density operator, $\Tr R=1$, and every fixed
$\alpha_b$ is finite. It is not a Gibbs state of the infinite cylinder.

\begin{proposition}[Admissibility of the complete preparation]\label{prop:state}
Each fixed Schwartz seminorm of $f_\lambda$ and each fixed mixed
parameter/spatial derivative is bounded on compact real $\lambda$ intervals.
For $\rho^I_\lambda=|f_\lambda\rangle\langle f_\lambda|\otimes R$ and every
fixed integer $p\ge0$, the quadratic-form sandwich is positive trace class,
with the exact payment
\begin{equation}\label{eq:fullmoment}
 \Tr(L_0^p\rho^I_\lambda L_0^p)
 =\sum_{j=0}^{2p}\binom{2p}{j}4^{2p-j}\alpha_{2p-j}
                                  \norm{P^jf_\lambda}_2^2<\infty.
\end{equation}
\end{proposition}
\begin{proof}
The matching exponents in \eqref{eq:sourcebounds} give
$|\Phi^{(\ell)}/\Phi|\le(B_\ell/\kappa)x^\ell$.
Set $a_\ell=\Phi^{(\ell)}/\Phi$ and $h=\sqrt\Phi$.
The identities $a_\ell'=a_{\ell+1}-a_\ell a_1$ yield the finite recurrence
\begin{equation}\label{eq:sqrtrecurrence}
 S_0=1,\qquad
 S_{j+1}=\tfrac12a_1S_j+
 \sum_{\ell=1}^j\frac{\partial S_j}{\partial a_\ell}
                          (a_{\ell+1}-a_\ell a_1),\qquad h^{(j)}=hS_j.
\end{equation}
Every monomial of $S_j$ has weighted degree $j$ if $a_\ell$ has weight
$\ell$. Hence
$|h^{(j)}|\le c_j\ee^{r/2}x^{j+1}\ee^{-x/2}$ for finite $c_j$.
The weakened exponent $-x/2$ still dominates every fixed polynomial,
Gaussian and logarithmic loss. Finite Leibniz differentiation therefore
bounds all spatial derivatives of $\ee^{\lambda v^2/2}h$ in Schwartz
seminorms. Parameter derivatives introduce only powers of $v^2/2$.
The integrals $Z_\lambda^{(a)}=\int v^{2a}F_\lambda\dd v$ and the positive
lower bound for $Z_\lambda$ control derivatives of $Z_\lambda^{-1/2}$.
Smooth positivity and complete reflection handle $v=0$.

For \eqref{eq:fullmoment}, expand $(k^2+4n^2)^{2p}$ in the full spectral
integral. Tonelli, Plancherel and \eqref{eq:R} give exactly the displayed
finite sum. This is the convergent positive sum of weighted ensemble
rank-one vectors; it does not assert that every vector in the density's
closed support belongs to $\Dom L_0^p$.
\end{proof}

All physical angular tails are the corresponding Gaussian series. For
fixed $j,N\ge0$ and $K>0$ the radial tail obeys
\begin{equation}\label{eq:radialtail}
 \int_{|k|>K}|k|^{2j}|\widehat f_\lambda(k)|^2\dd k
 \le K^{-2N}\norm{P^{j+N}f_\lambda}_2^2.
\end{equation}
Thus approximation pays the full spectra. Constants depend on the orders;
no order-uniform estimate follows. A positive Schwartz function does not in
general have a Schwartz square root: the ratio argument above is essential.
These are the preparation results of \cite[Section 3]{TNDDirect}, now in the
notation used for the count and collision calculus.

\section{A complete count instrument along the free motion}\label{sec:instrument}

\subsection{The transported quadrature}
Fix $\gamma>0$, set $T=\gamma^{-1}$, and keep
\begin{equation}\label{eq:Ut}
 U(t)=\ee^{-\ii\omega tL_0},\qquad Q_t=U(t)XU(t)^*.
\end{equation}
The space $\mathscr C=\Sch(\R)\otimes\K_{\rm fin}$, with finite sums
understood, is an invariant core for the required constituent powers.
Conjugating $\ii\partial_k$ by the Fourier multiplier
$\ee^{-\ii\omega t(k^2+4n^2)}$ gives
\begin{equation}\label{eq:Qt}
 Q_t=X-2\omega tP\quad\hbox{on }\mathscr C,\qquad
 \Dom Q_t^q=U(t)\Dom X^q.
\end{equation}
Unitary conjugation specifies the closed realization; no equality with
$\Dom X\cap\Dom P$ is required. Define the strong-law domain
\begin{equation}\label{eq:Dt}
 \mathcal D_t=\Dom L_0\cap\Dom Q_t^2
             =U(t)(\Dom L_0\cap\Dom X^2).
\end{equation}
Both the controller and readout follow $Q_t$. A fixed laboratory-position
attenuation would interact differently with the kinetic dispersion, since
$P^2$ and $X^2$ do not commute.

\subsection{Bounded outcomes, completeness and composition}
For $0\le s\le t\le T$ and $a=\gamma(t-s)$ define
\begin{equation}\label{eq:K}
 B_j(a)=\ee^{-aX^2/2}\frac{(\sqrt aX)^j}{\sqrt{j!}},\qquad
 K_j(t,s)=U(t)B_j(a)U(s)^*,\qquad j=0,1,2,\ldots.
\end{equation}
At $a=0$ put $B_0(0)=I$ and $B_j(0)=0$ for $j>0$. The signed factor
$X^j$ is retained, including for odd counts. The count label is an apparatus
outcome, not an additional angular multiplicity.

\begin{theorem}[Complete ideal instrument]\label{thm:instrument}
Each $K_j(t,s)$ is bounded. For a positive trace-class state $\rho$, the
operations
\[
 \mathcal I_{t,s}(J)(\rho)=\sum_{j\in J}K_j(t,s)\rho K_j(t,s)^*,
 \qquad J\subseteq\mathbb Z_{\ge0},
\]
converge in trace norm and form a countably additive completely positive
instrument. The all-count channel preserves trace. Sequential intervals
aggregate counts exactly, and the all-count channels compose.
\end{theorem}
\begin{proof}
For $a>0$, the scalar squared multiplier is $\ee^{-y}y^j/j!$, $y=av^2$.
Its maximum gives
\begin{equation}\label{eq:Knorm}
 \norm{B_0(a)}=1,\qquad
 \norm{B_j(a)}^2=\frac{j^j\ee^{-j}}{j!}\le1\quad(j\ge1).
\end{equation}
Bounded pointwise convergence proves the strong zero-duration limits.
Operator-norm continuity at zero is neither used nor generally true.
Let $c_N(y)=\ee^{-y}\sum_{j=0}^Ny^j/j!$. The positive exponential series
satisfies $0\le c_N\uparrow1$. For $0\le y\le R_0$ and $N+2\ge2R_0$,
its tail is bounded by $2R_0^{N+1}/(N+1)!$, by a geometric bound on the
successive ratios. On the complementary spectral interval it is at most
one. Consequently
\begin{equation}\label{eq:completeness}
 \sum_{j=0}^NK_j(t,s)^*K_j(t,s)
 =U(s)c_N(aX^2)U(s)^*\xrightarrow{\rm strong}I.
\end{equation}
Spectral decomposition of a positive trace-class $\rho$ and domination by
its summable eigenvalues show that the omitted input trace tends to zero.
Output tails are positive, so their trace norms equal their traces.
This proves trace-norm convergence and trace preservation. The map
$V\psi=(K_j\psi)_{j\ge0}$ is an isometry by \eqref{eq:completeness};
projection onto a subset of count coordinates followed by the count-space
partial trace proves complete positivity and countable additivity, including
finite matrix amplifications.

For $r\le s\le t$, set $a=\gamma(t-s)$, $b=\gamma(s-r)$ and $c=a+b>0$.
Direct multiplication gives, with $J=j+\ell$,
\begin{equation}\label{eq:composition}
 K_j(t,s)K_\ell(s,r)
 =\sqrt{\binom Jj}\left(\frac ac\right)^{j/2}
                     \left(\frac bc\right)^{\ell/2}K_J(t,r).
\end{equation}
The \emph{squared} coefficients sum to one for each $J$. Summing the
corresponding operations, rather than the amplitudes, therefore recovers
the aggregate-count operation. Positive trace-norm convergence justifies
regrouping. At zero durations the stated identity/zero convention applies.
Summing over $J$ gives composition of the all-count channels.
\end{proof}

The instrument is defined by bounded operations on all states. This proof
supplies the required infinite-count result directly; it does not invoke an
unbounded-generator master-equation theorem or a microscopic dilation.

\subsection{The no-count strong law}
\begin{theorem}[No-count propagation]\label{thm:strong}
The no-count operators satisfy $K_0(t,s)K_0(s,r)=K_0(t,r)$, send
$\mathcal D_s$ into $\mathcal D_t$, and obey
\begin{equation}\label{eq:strong}
 \partial_tK_0(t,s)\psi=
 \left[-\ii\omega L_0-\frac\gamma2Q_t^2\right]K_0(t,s)\psi,
 \qquad\psi\in\mathcal D_s.
\end{equation}
For fixed $s<T$, derivatives are right-sided at $t=s$ and left-sided at
$t=T$. The singleton $s=t=T$ has no nontrivial fixed-$s$ derivative
inside the prescribed time interval.
\end{theorem}
\begin{proof}
The cocycle follows from \eqref{eq:composition} at count zero. Put
$\chi=U(s)^*\psi\in\Dom L_0\cap\Dom X^2$ and
$m_a(v)=\ee^{-av^2/2}$. Its first two derivatives have the bounds
\[
 \norm{m_a'}_\infty\le\sqrt{a/\ee},\qquad
 \norm{m_a''}_\infty\le(1+2/\ee)a.
\]
The weak Sobolev product rule gives
\begin{equation}\label{eq:product}
 L_0(m_a\chi)=m_aL_0\chi-2m_a'\partial_v\chi-m_a''\chi,
 \qquad X^2m_a\chi=m_aX^2\chi.
\end{equation}
Since $\chi\in\Dom P^2\cap\Dom A$, all terms belong to $\HH$.
The angular operator commutes with $m_a$. These formulas prove preservation
of both constituent domains and graph continuity, including at $a=0$.
No mixed $XP$ assumption is inserted.
For admissible $a,a+h\ge0$, the mean-value bound
$|(m_{a+h}-m_a)/h|\le v^2/2$ gives
$\partial_a(m_a\chi)=-X^2m_a\chi/2$ strongly on $\Dom X^2$.
The free-unitary derivative holds on $\Dom L_0$ by full spectral calculus.
Combining their difference quotients and using graph continuity proves
\eqref{eq:strong}, together with the conjugated-domain assertion.
The same argument with one-sided quotients handles the endpoints.
\end{proof}

The effective expression is
$H_{\rm kin}-\ii\hbar\gamma Q_t^2/2$, with coupling notation
$\sqrt\gamma Q_t$. Its sign and half factor agree with the standard
no-count convention \cite[Section 1, (4)--(11)]{CastinDalibardMolmer2008v1}.
It is not asserted to be self-adjoint, and the theorem makes no claim about
a larger maximal domain for that formal effective generator. The domain
and convergence statements are exactly what the subsequent preparation
and observation require \cite[Section 4]{TNDDirect}.

\section{Conditioning and absolute Newman observation}\label{sec:observation}

Load the source once, at parameter $1/2$:
\begin{equation}\label{eq:initial}
 \rho_*=|f_{1/2}\rangle\langle f_{1/2}|\otimes R,\qquad
 \lambda(t)=\tfrac12-\gamma t,\qquad 0\le t\le T.
\end{equation}
The source-position map $u=v=x_{\rm phys}/\ell_0$ and this loader remain
supplied. The prepared vectors lie in every needed finite constituent
domain by Proposition~\ref{prop:state}.

\begin{theorem}[Exact conditional family and calibrated readout]\label{thm:readout}
Let $\widetilde\rho_t=K_0(t,0)\rho_*K_0(t,0)^*$ and
$p_t=\Tr\widetilde\rho_t$. Then
\begin{equation}\label{eq:conditional}
 \widetilde\rho_t=p_tU(t)\rho^I_{\lambda(t)}U(t)^*,\qquad
 p_t=\frac{Z_{\lambda(t)}}{Z_{1/2}}\in(0,1],\qquad
 \rho_t:=\frac{\widetilde\rho_t}{p_t}=U(t)\rho^I_{\lambda(t)}U(t)^*.
\end{equation}
For real $k$ the bounded effects
$E_\pm(t,k)=[I\pm\cos(kQ_t)]/2$ give
\begin{equation}\label{eq:q}
 q_\pm(t,k)=\Tr(\widetilde\rho_tE_\pm(t,k))
 =\frac12\left[p_t\pm\frac{H_{\lambda(t)}(k)}{H_{1/2}(0)}\right].
\end{equation}
Together with $q_{\rm fail}=1-p_t$, these are a normalized probability
law, and \eqref{eq:mainintro} holds on the closed time interval.
\end{theorem}
\begin{proof}
In the interaction frame,
\[
 \ee^{-\gamma tX^2/2}f_{1/2}
 =\sqrt{Z_{\lambda(t)}/Z_{1/2}}\,f_{\lambda(t)}.
\]
This proves \eqref{eq:conditional}. Positivity and compact-slab normalization
give $p_t\ge Z_{-1/2}/Z_{1/2}>0$; the attenuation is a contraction, so
$p_t\le1$. Real spectral calculus gives $0\le E_\pm\le I$ and
$E_++E_-=I$. Bounded-factor trace cyclicity and the complete angular
trace $\Tr R=1$ yield
\begin{equation}\label{eq:cosine}
 C(t,k):=\Tr(\rho_t\cos(kQ_t))
 =\int_\R\frac{F_{\lambda(t)}(v)}{Z_{\lambda(t)}}\cos(kv)\dd v
 =\frac{H_{\lambda(t)}(k)}{H_{\lambda(t)}(0)}.
\end{equation}
Multiplying by $p_t$ proves \eqref{eq:q}. On the initial state, the three
effects are $K_0^*E_+K_0$, $K_0^*E_-K_0$ and $I-K_0^*K_0$; they sum to
$I$ and are positive. Thus the quantities are actual joint probabilities
of the stipulated instrument and no-count readout, with the standard
operation normalization \cite[Section III, (17)--(19)]{JacobsSteck2006v1}.
\end{proof}

The success frequency is essential to absolute calibration. Selected
successful samples give $H_\lambda(k)/H_\lambda(0)$; multiplying by $p_t$
and the single initial scale gives $H_\lambda(k)$. At $k=0$,
$q_+=p_t$ and $q_-=0$. At $t=1/(2\gamma)$ the readout is precisely
\eqref{eq:xi}. These checks retain the full-line factor $Z_\lambda=2H_\lambda(0)$.
The observed density is the distribution of $Q_t$, not generally the
laboratory position $X$ after free dispersion.

\subsection{Regular state motion and signed evolution}
Write $w_t=|f_{\lambda(t)}|^2$ and $\mu_r(t)=\int v^r w_t(v)\dd v$.
Differentiation of the paid source integrals gives
\begin{align}
 \partial_tw_t&=-\gamma(v^2-\mu_2)w_t,
 &\dot p_t&=-\gamma\mu_2p_t,\label{eq:selection}\\
 \partial_tf_{\lambda(t)}&=-\frac\gamma2(X^2-\mu_2)f_{\lambda(t)},
 &\dot\mu_2&=-\gamma(\mu_4-\mu_2^2).\label{eq:tangent}
\end{align}
The first equation in \eqref{eq:tangent} includes both the square-root
half factor and the normalization subtraction. In particular
\begin{equation}\label{eq:speed}
 \norm{\partial_tf_{\lambda(t)}}_2^2
 =\frac{\gamma^2}{4}(\mu_4-\mu_2^2)<\infty.
\end{equation}
All fixed moments are bounded on the slab. The complete prepared density
operator is smooth to every fixed finite order in trace norm: differentiate
the rank-one ensemble and use the corresponding radial norms and summable
angular graph moments. This does not assert norm differentiability of
$U(t)$ on arbitrary states.

For a fixed polynomially bounded test function used here, the same
calculation gives
\begin{equation}\label{eq:covariance}
 \frac{d}{dt}\E_t a(X)=-\gamma\Cov_t(a(X),X^2).
\end{equation}
Consequently, with $d_0(t,k)=p_tC(t,k)$,
\begin{equation}\label{eq:heat}
 C_t=\gamma(C_{kk}+\mu_2C),\qquad
 (d_0)_t=\gamma(d_0)_{kk},\qquad
 d_0=\frac{H_{\lambda(t)}(k)}{H_{1/2}(0)}.
\end{equation}
The normalization term is absent only in the absolute contrast $d_0$.
The physical-time heat sign is positive because $\lambda(t)$ decreases.
The Fourier amplitude of $f_\lambda$, its momentum probability, and the
Fourier transform of $|f_\lambda|^2$ are different objects; only the last
produces this density-cosine observation. Entire continuation of
\eqref{eq:Hmixed} supplies no bounded nonreal-$k$ measurement.

\section{Signed moments in the complete count record}\label{sec:count}

The realization has more structure than its selected no-count output.
For each $j\ge0$, define
\begin{equation}\label{eq:countdef}
 \sigma_j(t)=K_j(t,0)\rho_*K_j(t,0)^*,\qquad
 P_j(t)=\Tr\sigma_j(t),\qquad
 d_j(t,k)=\Tr[\sigma_j(t)\cos(kQ_t)].
\end{equation}
Thus $P_0=p_t$ and $d_0$ is the absolute contrast in \eqref{eq:heat}.
For nonzero counts these are well-defined analytical traces of the existing
output operators. The frozen physical contract supplies the selected
no-count readout; mathematical definition of $d_j$ does not add a guarantee
of access to the same measurement after every nonzero count.

The literal count integral and normalization below were specifically
checked in \cite{TNDHighCountReview}. Section~\ref{sec:highcount} gives an
independent assessed complete-source proof of count summation and its first
real-$k$ derivative \cite{TNDCountTailReview}. The finite transfer laws and
whole-channel kernel calculation below retain their proof-author
diagnostic status; these limited assessments do not cover every claim of
\cite{TNDDirectCollision}.

\begin{theorem}[Count-sector moment and generating identities]\label{thm:counts}
For fixed $j\ge0$, $0\le t\le T$ and real $k$,
\begin{align}
 P_j(t)&=p_t\frac{(\gamma t)^j}{j!}\mu_{2j}(t),\label{eq:Pj}\\
 d_j(t,k)&=\frac{(\gamma t)^j}{j!Z_{1/2}}
      \int_\R v^{2j}F_{\lambda(t)}(v)\cos(kv)\dd v
       =\frac{(-\gamma t)^j}{j!}\partial_k^{2j}d_0(t,k),\label{eq:dj}\\
 |d_j(t,k)|&\le P_j(t),\qquad\sum_{j\ge0}P_j(t)=1.\label{eq:countabs}
\end{align}
At $t=0$ the stated count-zero convention applies. For $0\le\eta\le1$,
\begin{equation}\label{eq:generate}
 \sum_{j\ge0}\eta^jd_j(t,k)
 =\frac{H_{\lambda(t)+\eta\gamma t}(k)}{H_{1/2}(0)}.
\end{equation}
The series converges absolutely. In particular,
\begin{equation}\label{eq:conserved}
 \sum_{j\ge0}d_j(t,k)=C(0,k).
\end{equation}
\end{theorem}
\begin{proof}
In the interaction frame the radial output vector is
$\ee^{-\gamma tv^2/2}(\sqrt{\gamma t}\,v)^jf_{1/2}(v)/\sqrt{j!}$.
The signed odd-count amplitude is kept; its two occurrences in the diagonal
trace give $v^{2j}$. The complete angular trace is exactly one. This proves
the first integral formulas, including \eqref{eq:Pj}. Taking $2j$ derivatives
in $k$ is justified by \eqref{eq:Hmixed} and produces $(-1)^jv^{2j}$,
proving \eqref{eq:dj}. Positivity of $\sigma_j$ and the contraction norm of
the real cosine give \eqref{eq:countabs}; count completeness gives the sum.

Inside the source integral, the absolute sum of the terms weighted by
$\eta^j$ is dominated by the positive sum with the cosine replaced by one.
It is at most the initial probability mass for $0\le\eta\le1$.
The exponential series therefore gives
\[
 \frac1{Z_{1/2}}\int_\R F_{\lambda(t)}(v)
        \ee^{\eta\gamma tv^2}\cos(kv)\dd v
 =\frac{2H_{\lambda(t)+\eta\gamma t}(k)}{Z_{1/2}},
\]
which is \eqref{eq:generate}. At $\eta=1$ the parameter returns to $1/2$,
proving \eqref{eq:conserved}. No derivative-order summation is inferred
from fixed-order estimates: this series is justified by its own absolute
probability bound.
\end{proof}

Equation~\eqref{eq:Pj} is also a mixture of Poisson count laws against the
initial density, with scalar intensity $\gamma tv^2$. For $t>0$, every
$P_j>0$ and its normalized signed trace is
\begin{equation}\label{eq:conditionalj}
 \frac{d_j}{P_j}=\frac{(-1)^j\partial_k^{2j}C}{\mu_{2j}}.
\end{equation}
The count record therefore organizes the same complete weighted moments
that direct source differentiation supplies.

\begin{proposition}[Finite-sector transfer laws]\label{prop:transfer}
For fixed $j$ and $t>0$,
\begin{equation}\label{eq:transfer}
 \partial_td_j=\frac jt d_j-\frac{j+1}{t}d_{j+1},\qquad
 \partial_tP_j=\frac jt P_j-\frac{j+1}{t}P_{j+1}.
\end{equation}
Let $q_{j,\pm}=(P_j\pm d_j)/2$. These nonnegative analytical effect
probabilities obey the same equations, and in particular
\begin{equation}\label{eq:loss}
 \partial_tq_{0,\pm}=-q_{1,\pm}/t\le0,\qquad
 \partial_td_0=-d_1/t.
\end{equation}
\end{proposition}
\begin{proof}
Differentiate the integral in \eqref{eq:dj}. The prefactor contributes
$jd_j/t$; differentiating $F_{\lambda(t)}$ contributes $-\gamma v^2$,
whose integral is $(j+1)d_{j+1}/t$ with the stated sign. The same calculation
without the cosine proves the formula for $P_j$. Finite moments justify
each fixed derivative. Taking sum and difference gives the equations for
$q_{j,\pm}$. No infinite-series derivative is used.
\end{proof}

The quotients in \eqref{eq:transfer} have the required limits at zero from
the source formulas; they should not be evaluated as $0/0$. For example,
$d_1(t,k)/t\to\gamma\E_0[X^2\cos(kX)]$. Every statement concerns the
fixed slab and fixed positive $\gamma$, without an infinite-time or uniform
small-$\gamma$ conclusion.

\subsection{What is conserved when all counts are retained}
The conservation in \eqref{eq:conserved} concerns the co-moving cosine
observation. It does not make the full state stationary. The distinction
can be seen directly from the same Kraus amplitudes. In the interaction
frame, the all-count output of the prescribed initial state has the radial
operator-valued kernel
\begin{equation}\label{eq:dephase}
 f_{1/2}(v)f_{1/2}(w)\,\ee^{-\gamma t(v-w)^2/2}R.
\end{equation}
Indeed summing the kernel multipliers gives
\[
 \ee^{-\gamma t(v^2+w^2)/2}
       \sum_{j\ge0}\frac{(\gamma tvw)^j}{j!}
 =\ee^{-\gamma t(v-w)^2/2}.
\]
The absolute series is bounded by
$\exp[-\gamma t(|v|-|w|)^2/2]\le1$, so the kernels converge in
Hilbert--Schmidt norm by domination by $|f_{1/2}(v)f_{1/2}(w)|$.
Theorem~\ref{thm:instrument} already gives trace-norm convergence to the
same operator. The diagonal density is unchanged, while off-diagonal
coherence is attenuated. This kernel identity is an algebraic consequence
of the fixed instrument, not an additional physical evolution assumption.
It also explains why replacing signed odd-count amplitudes by their
absolute values would change the complete operation.

\section{Collisions in a regular conditional dynamics}\label{sec:collisions}

A real collision is a multiple zero of the signed readout, not a meeting
of eigenvalues of the fixed kinetic Hamiltonian. The state itself is
independent of which real observation parameter $k$ is queried.
Since $C(t,0)=1$, an event has $k_*\ne0$; evenness lets us discuss
$k_*>0$. The results of this section are conditional local descriptions,
not assertions that the canonical family has a collision.

\subsection{The probability and moment signature of a double event}
\begin{proposition}[Regular state, degenerate observation]\label{prop:event}
Suppose $0<t_*<T$ and
\begin{equation}\label{eq:double}
 C(t_*,k_*)=C_k(t_*,k_*)=0,\qquad A_*:=C_{kk}(t_*,k_*)\ne0.
\end{equation}
Then the success mass, state and count probabilities remain smooth, with
\begin{align}
 q_{0,+}=q_{0,-}&=p_*/2>0,&
 \partial_kq_{0,+}=\partial_kq_{0,-}&=0,\label{eq:eventprob}\\
 \partial_t(q_{0,+}-q_{0,-})&=\gamma p_*A_*\ne0,&
 d_1(t_*,k_*)&=-\gamma t_*p_*A_*\ne0.
 \label{eq:eventrate}
\end{align}
The event is equivalently described by the signed moments
\begin{equation}\label{eq:momenttriple}
 \E_*\cos(k_*X)=0,\qquad
 \E_*X\sin(k_*X)=0,\qquad
 \E_*X^2\cos(k_*X)=-A_*\ne0.
\end{equation}
\end{proposition}
\begin{proof}
Differentiate \eqref{eq:cosine} in $k$ using the finite moments.
This gives \eqref{eq:momenttriple}. At the event,
$(d_0)_{kk}=p_*A_*$ because $p_t$ is independent of $k$.
Equations~\eqref{eq:q}, \eqref{eq:heat} and \eqref{eq:dj} prove
\eqref{eq:eventprob}--\eqref{eq:eventrate}. Smoothness and the positive
floor for $p_t$ follow from Proposition~\ref{prop:state} and
\eqref{eq:selection}--\eqref{eq:speed}. The fixed-$j$ integrals give
smooth count probabilities. No division by the vanishing contrast is
needed for any state or probability formula.
\end{proof}

Both joint no-count outcome probabilities decrease by \eqref{eq:loss}.
At the event their values are equal, while their rates of decrease differ.
The one-count signed trace is exactly that difference multiplied by $-t_*$.
This is the dynamical content of the double-event curvature. It does not
make the event forbidden: the signed second moment has no sign constraint
from state positivity alone. Nor is $\mu_2$ a kinetic energy; its monotonicity
in \eqref{eq:tangent} is a property of the quadrature's second moment.

\subsection{Local zero geometry and the singularity of zero tracking}
\begin{theorem}[Local double-collision law]\label{thm:fold}
Under \eqref{eq:double}, let $a_*=(d_0)_{kk}(t_*,k_*)=p_*A_*$. Then
\begin{equation}\label{eq:fold}
 d_0(t_*+h,k_*+x)
 =\frac{a_*}{2}(x^2+2\gamma h)+O(|x|^3+|hx|+h^2).
\end{equation}
For $h<0$ sufficiently small there are two nearby real zeros,
\begin{equation}\label{eq:realroots}
 k_\pm(t)=k_*\pm\sqrt{-2\gamma h}+O(|h|).
\end{equation}
For $h>0$ the corresponding nearby analytic zeros form a conjugate pair,
\begin{equation}\label{eq:complexroots}
 z_\pm(t)=k_*\pm\ii\sqrt{2\gamma h}+O(|h|).
\end{equation}
All constants are local to the event. The complex roots refer only to
entire continuation of the source observation.
\end{theorem}
\begin{proof}
At the event $d_0=(d_0)_k=0$ and $(d_0)_t=\gamma a_*$ by
\eqref{eq:heat}. Finite multivariate Taylor expansion gives \eqref{eq:fold};
the required derivatives are paid by \eqref{eq:Hmixed}.
For $h=-s^2$ and $x=sy$, divide by $s^2$ and extend at $s=0$.
The limiting equation is $a_*(y^2/2-\gamma)=0$, with two simple real
roots. The analytic implicit-function theorem supplies real analytic
branches $y_\pm(s)=\pm\sqrt{2\gamma}+O(s)$, proving
\eqref{eq:realroots}. For $h=s^2$, the limiting roots are
$\pm\ii\sqrt{2\gamma}$, giving \eqref{eq:complexroots}.
A small fixed complex circle enclosing only the double zero at $h=0$
retains two zeros counted with multiplicity for small $h$, by uniform
analytic perturbation and the argument principle. The constructed roots
are therefore the full local pair. Reality of the source makes the
nonreal pair conjugate.
\end{proof}

Larger Newman parameter lies on the real-pair side because
$\lambda(t)-\lambda(t_*)=-\gamma h$. Reversing physical time reverses
the local picture. On a simple real zero curve $r(t)$,
\begin{equation}\label{eq:velocity}
 \dot r(t)=-\gamma\frac{(d_0)_{kk}}{(d_0)_k}
          =\frac{d_1}{t(d_0)_k}\quad(t>0).
\end{equation}
This follows by differentiating $d_0(t,r(t))=0$. Its divergence near the
double event belongs to locating the zero, while the state derivatives
remain finite. The squared gap has leading behavior
$(k_+-k_-)^2=-8\gamma h+o(|h|)$. Changing $\gamma$ rescales the clock;
it does not change an event's Newman parameter. The angular energies and
$\omega$ cancel from the co-moving observation, so they supply no additional
repulsion term in this local law.

For the original ternary observation $Y=+1,-1,0$ (the two successful
outcomes and failure),
\begin{equation}\label{eq:variance}
 \E Y=d_0,\qquad\Var Y=p_t-d_0^2.
\end{equation}
At the event this is $p_*>0$. A vanishing signed mean therefore has no
automatically vanishing single-trial noise. This exact variance calculation
is not a numerical experiment or a claim of optimal sample complexity.

\subsection{Balance of the positive and negative cosine sectors}
For fixed real $k\ne0$, let
$c_\pm(v)=\max(\pm\cos(kv),0)$ and define
\begin{equation}\label{eq:sectors}
 M_\pm=\E_t c_\pm(X)>0,\qquad
 B_\pm=\frac{\E_t[X^2c_\pm(X)]}{M_\pm},\qquad C=M_+-M_-.
\end{equation}
These are source-sector decompositions, distinct from the physical effects
$(I\pm\cos(kQ_t))/2$. Their boundaries are fixed when differentiating in
$t$, and finite moments justify the calculation despite the positive-part
corners. Equation~\eqref{eq:covariance} gives
\begin{equation}\label{eq:sectorratio}
 \partial_t\log(M_+/M_-)=-\gamma(B_+-B_-).
\end{equation}
At $C=0$, write $M_+=M_-=M$. Directly from the source integral,
\begin{equation}\label{eq:sectorcurvature}
 C_{kk}=-M(B_+-B_-),\qquad C_t=-\gamma M(B_+-B_-).
\end{equation}
Thus equal cosine masses, the stationary sine balance in
\eqref{eq:momenttriple}, and unequal sector second moments describe the
same double event. Quadratic attenuation changes the balance by suppressing
different radial locations at different rates. This interpretation and the
count-sector identity use the same signed source moments.

Cauchy--Schwarz gives the generic speed estimate
\begin{equation}\label{eq:covbound}
 |C_t|\le\gamma\sqrt{(\mu_4-\mu_2^2)
                   (\E_t\cos^2(kX)-C^2)}.
\end{equation}
It controls neither the sign nor the simultaneous vanishing of $C,C_k$.
It is available by direct source analysis and supplies no additional
collision margin merely through its dynamical interpretation.

\subsection{Higher multiplicities and temporal endpoints}
\begin{proposition}[Finite-order event polynomial]\label{prop:multiple}
If $d_0(t_*,\cdot)$ has exact finite order $m\ge2$ at $k_*$ and
$b_m=\partial_k^m d_0(t_*,k_*)\ne0$, then, uniformly for $(h,x)$ in a
fixed compact set,
\begin{equation}\label{eq:multiple}
 d_0(t_*+\delta^2h,k_*+\delta x)
 =\frac{b_m\delta^m}{m!}
 \sum_{r=0}^{\lfloor m/2\rfloor}
 \frac{m!\gamma^rh^rx^{m-2r}}{r!(m-2r)!}+o(\delta^m).
\end{equation}
For an interior event with $t_*>0$, if $m=2r$ then $d_j=0$ for $j<r$
and $d_r\ne0$. If $m=2r+1\ge3$, then $d_j=0$ for $j\le r$ and
$\partial_kd_r\ne0$.
\end{proposition}
\begin{proof}
Repeated differentiation of \eqref{eq:heat} gives
$\partial_t^a\partial_k^b d_0=\gamma^a\partial_k^{2a+b}d_0$.
In a finite multivariate Taylor expansion, every term of parabolic
weight $2a+b<m$ vanishes. Terms of weight $m$ give exactly
\eqref{eq:multiple}; a sufficiently high finite total-order remainder is
$o(\delta^m)$. The count assertions follow from \eqref{eq:dj} and the
nonzero factor $(-\gamma t_*)^j/j!$.
\end{proof}

The constants depend on $m$ and on the local event. A first-count diagnostic
therefore does not cover arbitrary multiplicities, and genericity is no
proof that the canonical family has only double events. At $t=0$ or $t=T$
only the physical side inside the slab is licensed. The entire source can
describe a two-sided germ without extending the apparatus beyond that slab.
The local laws are the familiar finite heat-flow geometry expressed in this
realization \cite{LocalCollisionGeometry,TNDDirectCollision}; no new
universal zero-collision law is claimed.

\section{An exact colliding control and the scope of the mechanism}\label{sec:control}

The instrument is source independent. Its application to a source uses
positivity, compact-parameter square-root regularity and complete moment
bounds. Consequently the generic identities must be tested on a source
that meets those requirements and nevertheless collides. The following
recorded reflected control \cite[Section 4, (9)--(12)]{ReflectedControl}
provides such a test with an exact event inside the physical slab.

\begin{theorem}[A smooth admissible source with a double event]\label{thm:control}
Let
\begin{align}
 g(v)&=\ee^{-v^4}+\tfrac12\ee^{-(v-4)^4}
                     +\tfrac12\ee^{-(v+4)^4},\label{eq:gcontrol}\\
 \Phi_*(v)&=c\ee^{-v^2/4}g(v),\qquad
 c=\left(16\int_0^\infty\cosh(v)\ee^{-v^2/4}g(v)\dd v\right)^{-1}>0.
 \label{eq:controlsource}
\end{align}
This source admits the same complete angular preparation and count
instrument with its own source normalization. It has all the fixed
square-root and moment domains used above. Its transform has a nondegenerate
double zero at $\lambda_*=1/4$, $k_*=\pi/4$, corresponding to laboratory
time $t_*=1/(4\gamma)$.
\end{theorem}
\begin{proof}
The source is strictly positive, even and smooth. Each derivative of a
summand of $g$ is that summand times a polynomial. Positivity of the mixture
therefore bounds $|g^{(\ell)}/g|$ by a fixed polynomial in $|v|$.
The finite ratio recurrence \eqref{eq:sqrtrecurrence} bounds every derivative
of $\sqrt g$ by a polynomial times $\sqrt g$. For $|v|\ge8$,
$|v-b|\ge|v|/2$ for $b=0,\pm4$, so
\begin{equation}\label{eq:controltail}
 g(v)\le2\ee^{-v^4/16}.
\end{equation}
Thus $\ee^{(\lambda-1/4)v^2/2}\sqrt g$ and its fixed mixed derivatives
have uniformly bounded Schwartz seminorms on each compact real $\lambda$
interval. All polynomial and exponential moments of the source are finite.
The normalization has a positive compact-interval lower bound.
With the same $R$, \eqref{eq:fullmoment} pays the full radial and angular
domains. The instrument proofs are unchanged.

At $\lambda=1/4$, ordinary full-line Fourier translation, followed by
evenness, gives
\begin{equation}\label{eq:controlfactor}
 H^*_{1/4}(k)=cF(k)(1+\cos4k),\qquad
 F(k)=\int_0^\infty\ee^{-v^4}\cos(kv)\dd v.
\end{equation}
Let $M_r=\int_0^\infty v^r\ee^{-v^4}\dd v$.
Integration of $(v\ee^{-v^4})'$ gives $M_4=M_0/4$.
Cauchy--Schwarz implies $M_2^2\le M_0M_4$, hence $M_2\le M_0/2$.
The global inequality $\cos y\ge1-y^2/2$ now gives
\begin{equation}\label{eq:Fpositive}
 F(k)\ge M_0(1-k^2/4)>0\qquad(|k|<2).
\end{equation}
At $k_*=\pi/4$ the factor $1+\cos4k$ and its first derivative vanish,
while its second derivative is $16$. Therefore
\begin{equation}\label{eq:controlcollision}
 H^*_{1/4}(k_*)=(H^*_{1/4})'(k_*)=0,\qquad
 (H^*_{1/4})''(k_*)=16cF(\pi/4)>0.
\end{equation}
Finally $1/2-\gamma t_*=1/4$ gives the stated interior physical time.
Every step is analytical, including the positivity of the curvature.
\end{proof}

The control has a smooth state, positive success mass, all count outputs,
the covariance law, the finite-sector transfer laws and the complete signed
sum in \eqref{eq:conserved}. It satisfies the local collision description
precisely. It fails the canonical integer-square arithmetic expansion and
the canonical theta tail formula; its own super-Gaussian bounds pay the
generic preparation requirements. Hence the theorem excludes an argument
based solely on those generic dynamical properties on a domain containing
this witness. It does not exclude source-sensitive analysis of the unchanged
theta source.

\subsection{What the other controls test}
The wider repository control collection \cite{CollisionControls} separates
premises rather than providing one universal counterexample class:
\begin{center}\small
\begin{tabularx}{\textwidth}{@{}>{\raggedright\arraybackslash}p{0.25\textwidth}Y@{}}
\toprule
Control family & Relevant distinction \\\midrule
Euler-factor sources & Fail the unchanged canonical arithmetic source; applicable positive smooth preparation and Fourier/filtering identities can survive. \\
Weighted-square and finite-prime changes & Complete unit weights or full canonical reflection fail. Finite endpoint matches do not give the smooth full-line source domains. \\
Reflected smooth and positive-primitive sources & Exact theta arithmetic is absent. Positivity, reflection and the applicable generic transport laws need not prevent collisions. \\
Atomic or nonsmooth finite profiles & Fourier moments can exist while the present smooth square-root and graph-domain requirements fail or need a different construction. \\
\bottomrule
\end{tabularx}
\end{center}
Only Theorem~\ref{thm:control}'s preparation and event are checked afresh
here; this table does not claim a new domain audit of every family. A finite
collision-control estimate should be tested where the relevant witness lies,
not only in a different collision-free region.

The analytical gain obtained so far is a complete representation and a
precise collision dictionary. Equations~\eqref{eq:covariance},
\eqref{eq:dj}, \eqref{eq:generate} and \eqref{eq:sectorcurvature} are also
available by differentiating the same source integral. Their reorganization
has not yet supplied an additional signed margin. A new mechanism must
identify a consequence of the canonical source that survives its deformation
and distinguishes the relevant controls quantitatively or structurally.

\section{Outstanding analytical tasks}\label{sec:outstanding}

The construction realizes the canonical family within the supplied ideal
conditional model. The collision analysis describes a hypothetical event.
Section~\ref{sec:highcount} now adds two assessed scalar consequences:
complete high-count localization with both signs, and absolute value/first-
derivative errors for finite count reconstruction. These settle neither
the source-specific collision question nor an advantage over direct source
analysis. They define useful constraints on the next conceptual pass.

\subsection{Identify an independent source-specific consequence}
The event conditions are the two vanished balances in
\eqref{eq:momenttriple} and, at a double event, a nonzero signed second
moment. Source-sector covariance and count-sector transfers describe the
same signed information. Equation~\eqref{eq:scorebalance} is also a
consequence of the existing balances, not an additional nonvanishing
condition. Theorem~\ref{thm:highcount} excludes eventual one-sign literal
count sectors for theta itself at fixed nonzero real $k$; renaming that
premise cannot supply a collision mechanism. The absolute bounds in
Theorem~\ref{thm:counttail} pay omitted sectors but do not create a
relative margin when the readout and its first derivative vanish.

The next question is whether the complete source imposes an independently
proved joint relation or quantitative constraint on the signed balances,
or yields another useful analytical consequence. State the actual source
premise, parameter/multiplicity range and control that tests it. A rigidity
or transport relation can be useful before a scalar bound is selected,
but requiring $C$ and $C_k$ never to vanish together only restates the
target. All global proposals retain the endpoint, confinement and
equivalence qualifications of existing heat-flow criteria.

\subsection{Freeze a bounded task only after a mathematical decision}
TND-D-B02's selected diagnostic and the E03/E04 and E06/E07 proof/assessment
lanes are complete within their stated scopes. The next-session brainstorm
is recorded as TND-D-B03, recommended Astra at effort ultra, ready but
unactivated at this manuscript integration. No new analytical execution
contract is selected here.

A subsequent contract must name one theorem, identity or estimate, the
source information, parameter/readout ranges, derivative orders, complete-
tail method and success or failure test. Compare with direct use of
exactly the same source and uncertainty information. The count labels add
no independent arithmetic data. The generic covariance bound
\eqref{eq:covbound} controls speed without controlling the dangerous signed
balance, and \eqref{eq:variance} gives nonzero observation variance at zero
contrast. A conditioning or practical reconstruction claim therefore
needs an explicit information/error model. The present conservative
threshold is finite but uncomputed and unoptimized; it is not a cost claim.

\subsection{Cover multiplicities and the desired global range}
The double-event picture is local and nondegenerate.
Proposition~\ref{prop:multiple} identifies higher-order data but supplies
neither constants uniform in multiplicity nor a mechanism covering all
events. An exclusion must cover the relevant multiplicities or justify
why omitted ones cannot arise. It must also handle temporal endpoints and
large readout parameters. The compact-real-$k$ phase theorem and the
all-real-$k$ absolute tail theorem have different scopes; neither supplies
an all-height relative collision margin. Fixed finite-order domains do not
imply a uniform graph or energy estimate as the count increases. The
closed apparatus slab gives no infinite-time conclusion.

\subsection{Separate analytical reconstruction from physical access}
The theta loader, spatial imprint, complete angular ensemble, absolute
scale and transported no-count readout remain supplied. Deriving them
intrinsically or specifying a microscopic implementation and its resource
or error bounds is a separate physical task. The nonzero-count quantities
are valid analytical traces, and \eqref{eq:tailreconstruction} is an
analytical finite reconstruction. An experiment using those joint
readouts needs its own explicit access specification. A changed controller,
source, clock or measurement law requires a revised construction before
transferring a claim. Changing $\gamma$ only changes the laboratory clock.

\subsection{Preserve assessment boundaries and use computation as support}
The realization retains its internal conditional assessment. The new
concentration, phase, complete-tail and first-derivative reconstruction
results have separate-worker, same-model internal LLM assessments with
shared premises. Only the specifically consumed count identities and
source-score identity share that acceptance. Other fold, higher-
multiplicity, finite-transfer and whole-channel kernel diagnostics retain
proof-author status; manuscript production does not supply specialist
independence or a wider assessment.

No numerical evidence is needed for these analytical results. Supporting
computation may later evaluate a declared constant, compare conditioning
under a fixed error model or test a finite-domain question, with complete-
tail errors and a distinction between floating and certified evidence.
It cannot replace the independent signed source estimate needed for
collision control. No new numerical campaign is selected by this draft.

\section*{Acknowledgements and evidence provenance}
This manuscript was drafted for Edward Baker with substantial assistance
from OpenAI's \texttt{gpt-6-astra} through Codex, with recorded reasoning
effort \texttt{max}. The complete-realization proof and its internal
assessment were prepared with \texttt{gpt-6.1-sol}, effort \texttt{ultra};
the construction and collision reconnaissance involved Astra-assisted
research discussions. Assistance included mathematical exposition, source
comparison, consistency checks and document production. The realization's
existing acceptance is internal, with shared premises and same-assistant
review. The subsequent collision/count calculations retain proof-author
status. This manuscript's local editorial checks are not independent
specialist review or formal verification and do not upgrade those statuses.
\begingroup\small
\bibliographystyle{plainurl}
\bibliography{direct_newman_sources}
\endgroup
\end{document}
